\unchapter{Conclusions}

This Bachelor's thesis investigated the architectural trade-offs between a modular monolithic architecture (\texttt{eShopOnWeb}) and a distributed microservices architecture (\texttt{eShopOnContainers}) for small and medium-sized e-commerce enterprises. Based on controlled Docker and \texttt{k6} benchmarks, source-level change propagation analysis across five standardized evolution scenarios, and triangulation with industrial migration studies, the following conclusions are drawn:

\begin{enumerate}[leftmargin=1.5em]
  \item \textbf{Significant Memory Efficiency of the Modular Monolith:} Under SME-scale host constraints ($5.0\text{ GiB}$ RAM ceiling), the modular monolith consumes $25.5\%$ less memory at idle ($903.98\text{ MiB}$ vs.\ $1,213.60\text{ MiB}$) and $14.7\%$ less memory under $20\text{ VUs}$ concurrent load ($1,151.35\text{ MiB}$ vs.\ $1,350.20\text{ MiB}$), directly lowering baseline cloud hosting costs.
  \item \textbf{Superior Tail Latency in Microservices:} Although both architectures deliver comparable sustained throughput ($7.49\text{ req/s}$ vs.\ $7.58\text{ req/s}$) with $0.00\%$ HTTP errors, microservices achieve $2.06\times$ faster median latency ($8.08\text{ ms}$ vs.\ $16.67\text{ ms}$) and $2.65\times$ faster $p_{95}$ tail latency ($12.30\text{ ms}$ vs.\ $32.57\text{ ms}$) through decoupled read endpoints.
  \item \textbf{Asymmetric Change Propagation:} Microservices isolate single-domain modifications ($2$--$3$ files in a single deployable container) but amplify cross-cutting changes ($9$ files across $3$ deployment units for multi-currency support vs.\ $6$ files in $1$ deployment unit for the monolith).
  \item \textbf{Actionable SME Recommendation:} SMEs with engineering teams of six or fewer developers should adopt a \textbf{Modular-Monolith-First} architecture enforced by strict Domain-Driven Design module boundaries, transitioning to selective microservice extraction (beginning with read-heavy Catalog services) only when sustained concurrency bottlenecks or multi-team organizational growth justify the operational overhead.
\end{enumerate}
